% Options for packages loaded elsewhere
% Options for packages loaded elsewhere
\PassOptionsToPackage{unicode}{hyperref}
\PassOptionsToPackage{hyphens}{url}
\PassOptionsToPackage{dvipsnames,svgnames,x11names}{xcolor}
%
\documentclass[
  letterpaper,
  DIV=11,
  numbers=noendperiod]{scrartcl}
\usepackage{xcolor}
\usepackage{amsmath,amssymb}
\setcounter{secnumdepth}{-\maxdimen} % remove section numbering
\usepackage{iftex}
\ifPDFTeX
  \usepackage[T1]{fontenc}
  \usepackage[utf8]{inputenc}
  \usepackage{textcomp} % provide euro and other symbols
\else % if luatex or xetex
  \usepackage{unicode-math} % this also loads fontspec
  \defaultfontfeatures{Scale=MatchLowercase}
  \defaultfontfeatures[\rmfamily]{Ligatures=TeX,Scale=1}
\fi
\usepackage{lmodern}
\ifPDFTeX\else
  % xetex/luatex font selection
\fi
% Use upquote if available, for straight quotes in verbatim environments
\IfFileExists{upquote.sty}{\usepackage{upquote}}{}
\IfFileExists{microtype.sty}{% use microtype if available
  \usepackage[]{microtype}
  \UseMicrotypeSet[protrusion]{basicmath} % disable protrusion for tt fonts
}{}
\makeatletter
\@ifundefined{KOMAClassName}{% if non-KOMA class
  \IfFileExists{parskip.sty}{%
    \usepackage{parskip}
  }{% else
    \setlength{\parindent}{0pt}
    \setlength{\parskip}{6pt plus 2pt minus 1pt}}
}{% if KOMA class
  \KOMAoptions{parskip=half}}
\makeatother
% Make \paragraph and \subparagraph free-standing
\makeatletter
\ifx\paragraph\undefined\else
  \let\oldparagraph\paragraph
  \renewcommand{\paragraph}{
    \@ifstar
      \xxxParagraphStar
      \xxxParagraphNoStar
  }
  \newcommand{\xxxParagraphStar}[1]{\oldparagraph*{#1}\mbox{}}
  \newcommand{\xxxParagraphNoStar}[1]{\oldparagraph{#1}\mbox{}}
\fi
\ifx\subparagraph\undefined\else
  \let\oldsubparagraph\subparagraph
  \renewcommand{\subparagraph}{
    \@ifstar
      \xxxSubParagraphStar
      \xxxSubParagraphNoStar
  }
  \newcommand{\xxxSubParagraphStar}[1]{\oldsubparagraph*{#1}\mbox{}}
  \newcommand{\xxxSubParagraphNoStar}[1]{\oldsubparagraph{#1}\mbox{}}
\fi
\makeatother

\usepackage{color}
\usepackage{fancyvrb}
\newcommand{\VerbBar}{|}
\newcommand{\VERB}{\Verb[commandchars=\\\{\}]}
\DefineVerbatimEnvironment{Highlighting}{Verbatim}{commandchars=\\\{\}}
% Add ',fontsize=\small' for more characters per line
\usepackage{framed}
\definecolor{shadecolor}{RGB}{241,243,245}
\newenvironment{Shaded}{\begin{snugshade}}{\end{snugshade}}
\newcommand{\AlertTok}[1]{\textcolor[rgb]{0.68,0.00,0.00}{#1}}
\newcommand{\AnnotationTok}[1]{\textcolor[rgb]{0.37,0.37,0.37}{#1}}
\newcommand{\AttributeTok}[1]{\textcolor[rgb]{0.40,0.45,0.13}{#1}}
\newcommand{\BaseNTok}[1]{\textcolor[rgb]{0.68,0.00,0.00}{#1}}
\newcommand{\BuiltInTok}[1]{\textcolor[rgb]{0.00,0.23,0.31}{#1}}
\newcommand{\CharTok}[1]{\textcolor[rgb]{0.13,0.47,0.30}{#1}}
\newcommand{\CommentTok}[1]{\textcolor[rgb]{0.37,0.37,0.37}{#1}}
\newcommand{\CommentVarTok}[1]{\textcolor[rgb]{0.37,0.37,0.37}{\textit{#1}}}
\newcommand{\ConstantTok}[1]{\textcolor[rgb]{0.56,0.35,0.01}{#1}}
\newcommand{\ControlFlowTok}[1]{\textcolor[rgb]{0.00,0.23,0.31}{\textbf{#1}}}
\newcommand{\DataTypeTok}[1]{\textcolor[rgb]{0.68,0.00,0.00}{#1}}
\newcommand{\DecValTok}[1]{\textcolor[rgb]{0.68,0.00,0.00}{#1}}
\newcommand{\DocumentationTok}[1]{\textcolor[rgb]{0.37,0.37,0.37}{\textit{#1}}}
\newcommand{\ErrorTok}[1]{\textcolor[rgb]{0.68,0.00,0.00}{#1}}
\newcommand{\ExtensionTok}[1]{\textcolor[rgb]{0.00,0.23,0.31}{#1}}
\newcommand{\FloatTok}[1]{\textcolor[rgb]{0.68,0.00,0.00}{#1}}
\newcommand{\FunctionTok}[1]{\textcolor[rgb]{0.28,0.35,0.67}{#1}}
\newcommand{\ImportTok}[1]{\textcolor[rgb]{0.00,0.46,0.62}{#1}}
\newcommand{\InformationTok}[1]{\textcolor[rgb]{0.37,0.37,0.37}{#1}}
\newcommand{\KeywordTok}[1]{\textcolor[rgb]{0.00,0.23,0.31}{\textbf{#1}}}
\newcommand{\NormalTok}[1]{\textcolor[rgb]{0.00,0.23,0.31}{#1}}
\newcommand{\OperatorTok}[1]{\textcolor[rgb]{0.37,0.37,0.37}{#1}}
\newcommand{\OtherTok}[1]{\textcolor[rgb]{0.00,0.23,0.31}{#1}}
\newcommand{\PreprocessorTok}[1]{\textcolor[rgb]{0.68,0.00,0.00}{#1}}
\newcommand{\RegionMarkerTok}[1]{\textcolor[rgb]{0.00,0.23,0.31}{#1}}
\newcommand{\SpecialCharTok}[1]{\textcolor[rgb]{0.37,0.37,0.37}{#1}}
\newcommand{\SpecialStringTok}[1]{\textcolor[rgb]{0.13,0.47,0.30}{#1}}
\newcommand{\StringTok}[1]{\textcolor[rgb]{0.13,0.47,0.30}{#1}}
\newcommand{\VariableTok}[1]{\textcolor[rgb]{0.07,0.07,0.07}{#1}}
\newcommand{\VerbatimStringTok}[1]{\textcolor[rgb]{0.13,0.47,0.30}{#1}}
\newcommand{\WarningTok}[1]{\textcolor[rgb]{0.37,0.37,0.37}{\textit{#1}}}

\usepackage{longtable,booktabs,array}
\newcounter{none} % for unnumbered tables
\usepackage{calc} % for calculating minipage widths
% Correct order of tables after \paragraph or \subparagraph
\usepackage{etoolbox}
\makeatletter
\patchcmd\longtable{\par}{\if@noskipsec\mbox{}\fi\par}{}{}
\makeatother
% Allow footnotes in longtable head/foot
\IfFileExists{footnotehyper.sty}{\usepackage{footnotehyper}}{\usepackage{footnote}}
\makesavenoteenv{longtable}
\usepackage{graphicx}
\makeatletter
\newsavebox\pandoc@box
\newcommand*\pandocbounded[1]{% scales image to fit in text height/width
  \sbox\pandoc@box{#1}%
  \Gscale@div\@tempa{\textheight}{\dimexpr\ht\pandoc@box+\dp\pandoc@box\relax}%
  \Gscale@div\@tempb{\linewidth}{\wd\pandoc@box}%
  \ifdim\@tempb\p@<\@tempa\p@\let\@tempa\@tempb\fi% select the smaller of both
  \ifdim\@tempa\p@<\p@\scalebox{\@tempa}{\usebox\pandoc@box}%
  \else\usebox{\pandoc@box}%
  \fi%
}
% Set default figure placement to htbp
\def\fps@figure{htbp}
\makeatother


% definitions for citeproc citations
\NewDocumentCommand\citeproctext{}{}
\NewDocumentCommand\citeproc{mm}{%
  \begingroup\def\citeproctext{#2}\cite{#1}\endgroup}
\makeatletter
 % allow citations to break across lines
 \let\@cite@ofmt\@firstofone
 % avoid brackets around text for \cite:
 \def\@biblabel#1{}
 \def\@cite#1#2{{#1\if@tempswa , #2\fi}}
\makeatother
\newlength{\cslhangindent}
\setlength{\cslhangindent}{1.5em}
\newlength{\csllabelwidth}
\setlength{\csllabelwidth}{3em}
\newenvironment{CSLReferences}[2] % #1 hanging-indent, #2 entry-spacing
 {\begin{list}{}{%
  \setlength{\itemindent}{0pt}
  \setlength{\leftmargin}{0pt}
  \setlength{\parsep}{0pt}
  % turn on hanging indent if param 1 is 1
  \ifodd #1
   \setlength{\leftmargin}{\cslhangindent}
   \setlength{\itemindent}{-1\cslhangindent}
  \fi
  % set entry spacing
  \setlength{\itemsep}{#2\baselineskip}}}
 {\end{list}}
\usepackage{calc}
\newcommand{\CSLBlock}[1]{\hfill\break\parbox[t]{\linewidth}{\strut\ignorespaces#1\strut}}
\newcommand{\CSLLeftMargin}[1]{\parbox[t]{\csllabelwidth}{\strut#1\strut}}
\newcommand{\CSLRightInline}[1]{\parbox[t]{\linewidth - \csllabelwidth}{\strut#1\strut}}
\newcommand{\CSLIndent}[1]{\hspace{\cslhangindent}#1}



\setlength{\emergencystretch}{3em} % prevent overfull lines

\providecommand{\tightlist}{%
  \setlength{\itemsep}{0pt}\setlength{\parskip}{0pt}}



 


\KOMAoption{captions}{tableheading}
\makeatletter
\@ifpackageloaded{caption}{}{\usepackage{caption}}
\AtBeginDocument{%
\ifdefined\contentsname
  \renewcommand*\contentsname{Table of contents}
\else
  \newcommand\contentsname{Table of contents}
\fi
\ifdefined\listfigurename
  \renewcommand*\listfigurename{List of Figures}
\else
  \newcommand\listfigurename{List of Figures}
\fi
\ifdefined\listtablename
  \renewcommand*\listtablename{List of Tables}
\else
  \newcommand\listtablename{List of Tables}
\fi
\ifdefined\figurename
  \renewcommand*\figurename{Figure}
\else
  \newcommand\figurename{Figure}
\fi
\ifdefined\tablename
  \renewcommand*\tablename{Table}
\else
  \newcommand\tablename{Table}
\fi
}
\@ifpackageloaded{float}{}{\usepackage{float}}
\floatstyle{ruled}
\@ifundefined{c@chapter}{\newfloat{codelisting}{h}{lop}}{\newfloat{codelisting}{h}{lop}[chapter]}
\floatname{codelisting}{Listing}
\newcommand*\listoflistings{\listof{codelisting}{List of Listings}}
\makeatother
\makeatletter
\makeatother
\makeatletter
\@ifpackageloaded{caption}{}{\usepackage{caption}}
\@ifpackageloaded{subcaption}{}{\usepackage{subcaption}}
\makeatother
\usepackage{bookmark}
\IfFileExists{xurl.sty}{\usepackage{xurl}}{} % add URL line breaks if available
\urlstyle{same}
\hypersetup{
  pdftitle={quartobot: Quarto-native automation for citation resolution and reproducible publishing},
  pdfkeywords={Quarto, manubot, manuscript-as-software, persistent
identifiers, citation resolution, reproducible publishing, scholarly
communication, CSL-JSON, MCP, pre-render hooks},
  colorlinks=true,
  linkcolor={blue},
  filecolor={Maroon},
  citecolor={Blue},
  urlcolor={Blue},
  pdfcreator={LaTeX via pandoc}}


\title{quartobot: Quarto-native automation for citation resolution and
reproducible publishing}
\author{Sean Davis}
\date{}
\begin{document}
\maketitle
\begin{abstract}
TODO --- Sean's section. \textasciitilde250--300 words. Four moves:
problem (manubot pattern, Quarto's broader publishing surface,
pandoc-manubot-cite covers citation resolution but not the broader
pattern); approach (pre-render hook architecture, committed
\texttt{references.json} as a durable artifact, CLI + MCP + on-disk
surfaces); results (wraps \texttt{citekey\_to\_csl\_item}, installable
via \texttt{uv\ tool\ install}, validated on two production
manuscripts); conclusion (extends manuscript-as- software along five
dimensions while preserving manubot's resolver expertise). Full bullets
in outline.md.
\end{abstract}


\section{Introduction}\label{introduction}

\section{Implementation}\label{implementation}

quartobot is structured as a single Python package exposing a small
command-line interface. The integration with Quarto runs through a
\texttt{project.pre-render:} hook in \texttt{\_quarto.yml}; the package
is otherwise independent of Quarto and can be invoked directly. Citation
resolution itself is delegated to manubot's
\texttt{citekey\_to\_csl\_item} function, so the resolver chain
(Crossref for DOI, NCBI E-utilities for PubMed and PMC, the arXiv API,
Open Library for ISBN, the Wikidata API for Wikidata items) is the one
already validated by eight years of manubot use.

\subsection{Architecture and the choice of pre-render over
filter}\label{architecture-and-the-choice-of-pre-render-over-filter}

Quarto's documentation describes two principal customization surfaces:
pre-render scripts, which run before pandoc parses the document, and Lua
filters, which run on the parsed abstract syntax tree after the document
engine executes. Citation \textbf{formatting} --- how a resolved
reference appears once pandoc-citeproc has it in hand --- is a natural
Lua-filter concern, and indeed pandoc-citeproc is itself implemented as
a filter. Citation \textbf{resolution}, in contrast, must happen before
the AST exists, because the bibliography is consumed at parse time. A
Lua filter is therefore unable to add entries to the bibliography; by
the time it runs, that file has already been read.

quartobot lives on the pre-render side for this reason. The
\texttt{resolve} subcommand scans Quarto sources (\texttt{.qmd},
\texttt{.md}, and \texttt{.ipynb}) for citation keys of the form
\texttt{@\textless{}prefix\textgreater{}:\textless{}identifier\textgreater{}},
resolves each through \texttt{citekey\_to\_csl\_item}, and writes the
resulting CSL-JSON to a file named \texttt{references.json}.
Pandoc-citeproc then consumes that file as part of the normal render.
The integration touches three things in a project: a single line under
\texttt{project:} in \texttt{\_quarto.yml}, the generated
\texttt{references.json}, and optionally a hand-curated
\texttt{references.bib} for entries the resolver does not cover
(textbooks, gray literature, software releases with no DOI).

This choice differs from the existing path via the
\texttt{pandoc-manubot-cite} filter, which provides the same resolver
behavior as a pandoc filter and therefore runs during each pandoc
invocation. The two approaches are not equivalent. A filter resolves at
render time and produces no durable output. A pre-render hook resolves
once, writes a file, commits it to the repository, and lets every
subsequent render consume the cached bibliography without further
network activity. For manuscripts rendered repeatedly --- by continuous
integration, by per-pull-request preview, by per-commit permalink ---
this difference becomes material. The bibliography becomes a versioned
artifact in its own right.

The trade-off favors pre-render when several of the following hold: the
same manuscript is rendered repeatedly, authors want offline-capable
renders, the bibliography is treated as a versioned artifact reviewed
alongside the prose, and determinism matters more than always-fresh
metadata. These are precisely the conditions under which the
manuscript-as-software pattern operates. Pre-render is therefore the
structurally correct architectural choice for the pattern;
per-invocation filters remain a valid choice for one-off pandoc renders
where the manuscript is not under continuous integration.

\subsection{Supported identifiers}\label{supported-identifiers}

Citation keys take the form
\texttt{@\textless{}prefix\textgreater{}:\textless{}identifier\textgreater{}}
in prose. The supported prefixes and their upstream resolvers are listed
in Table 1.

{\def\LTcaptype{none} % do not increment counter
\begin{longtable}[]{@{}lll@{}}
\toprule\noalign{}
Prefix & Identifier type & Upstream resolver \\
\midrule\noalign{}
\endhead
\bottomrule\noalign{}
\endlastfoot
\texttt{@doi:} & Digital Object Identifier & Crossref \\
\texttt{@pmid:} & PubMed identifier & NCBI E-utilities \\
\texttt{@pmc:} & PubMed Central identifier & NCBI E-utilities \\
\texttt{@arxiv:} & arXiv identifier & arXiv API \\
\texttt{@isbn:} & International Standard Book Number & Open Library \\
\texttt{@wikidata:} & Wikidata QID & Wikidata API \\
\texttt{@url:} & Web resource URL & URL metadata extraction \\
\end{longtable}
}

Hand-curated entries continue to be supported through a project-level
\texttt{references.bib}, which pandoc reads alongside the resolver
output. The choice of which mechanism to use for a given citation is the
author's; the two coexist in the same document. Reconciliation between
the two sources is described below.

\subsection{Input format coverage}\label{input-format-coverage}

The pre-render scan is format-agnostic on the input side. Jupyter
notebooks (\texttt{.ipynb}), Quarto Markdown (\texttt{.qmd}), and plain
Markdown (\texttt{.md}) are all scanned in the same pass, because Quarto
presents these formats through a uniform interface. This is the
dimension manubot's bespoke toolchain did not reach: notebooks were
never first-class manuscript inputs. The same \texttt{@doi:} syntax that
resolves in prose resolves equally in a notebook narrative cell,
allowing a literate-programming manuscript (analysis code, narrative,
and citations in the same file) to use the manuscript-as-software
pattern without modification.

\subsection{Cross-platform deployment}\label{cross-platform-deployment}

quartobot inherits Quarto's platform reach. The package runs on macOS,
Linux, and Windows; the GitHub Actions workflows that ship with
\texttt{quartobot\ use\ github-ci} (described below) use standard
cross-platform actions. Manubot's bespoke toolchain assumed a
Linux+Docker rendering environment; quartobot makes no such assumption.

\subsection{CI scaffolding}\label{ci-scaffolding}

The \texttt{use\ github-ci} subcommand lays down the
continuous-integration scaffolding for the manuscript-as-software
pattern. The default configuration produces three workflows. The first
builds the manuscript on every push to the default branch, rendering to
HTML, PDF, DOCX, and JATS XML, and publishes the result to GitHub Pages.
The second builds a per-pull-request preview at a stable URL under
\texttt{…/pr/\textless{}n\textgreater{}/}, posting a sticky comment with
the preview link for reviewers who do not have a local Quarto
installation. The third generates a versions index on \texttt{gh-pages}
listing tagged releases and open pull-request previews.

With the optional \texttt{-\/-with-versioned-snapshots} flag, each
commit on the default branch also produces an immutable snapshot at
\texttt{…/v/\textless{}sha\textgreater{}/}. This is the manubot
per-commit permalink pattern, preserved exactly: a rendered manuscript
on disk always knows which commit produced it, and a downloaded copy
carries that provenance in the HTML.

\subsection{Three surfaces for human and AI
use}\label{three-surfaces-for-human-and-ai-use}

quartobot exposes its resolution functionality through three surfaces,
each appropriate for a different consumer.

The \textbf{CLI} (\texttt{quartobot\ resolve}) is the primary interface
and the one Quarto invokes through the pre-render hook. The same CLI is
the surface AI authoring agents shell out to when drafting. An agent
that has just written \texttt{@doi:10.1371/journal.pcbi.1007128} into
prose can pass that identifier, or a batch of several dozen, to a single
\texttt{quartobot\ resolve} call and receive the resolved CSL-JSON back.
Two checks become cheap. First, the agent can confirm that every
identifier resolved at all; an identifier that does not resolve is one
the agent likely fabricated, and the citation can be removed or
corrected before it reaches the manuscript. Second, the agent can
compare the resolved metadata to what it intended (the work's title, the
lead author's surname, the year of publication) and catch the case where
the identifier resolves but does not name the paper the agent meant to
cite. The pattern matches the broader recommendation in recent work on
agent-tool efficiency to expose tools as code on a filesystem rather
than as in-context schemas ({``Code Execution with {MCP:} Building More
Efficient {AI} Agents \textbackslash{} {Anthropic},''} n.d.).

The \textbf{MCP server} (\texttt{quartobot\ mcp}) is an optional
stdio-based Model Context Protocol surface for the case where an IDE or
agent framework already operates inside an MCP-aware environment and
shelling out adds ergonomic friction. The server exposes the same
resolver function as the CLI; results are identical.

The \textbf{on-disk artifact} (\texttt{references.json}) is the surface
that distinguishes quartobot most sharply from a filter-based approach.
Once \texttt{resolve} has written the file and an author has committed
it, the resolved citation graph is available to any downstream consumer
without further network activity and without invoking pandoc. Citation
linters can flag unused references, duplicate identifiers, or missing
required fields. Bibliometric pipelines can ingest a manuscript's
reference set as structured data. Other rendering toolchains can read
the JSON directly. AI authoring agents drafting later prose can ground
new claims against the manuscript's existing reference set, reading the
committed file rather than re-querying the upstream APIs. This reuse is
what \texttt{pandoc-manubot-cite} cannot offer: filter outputs are
produced and consumed within a single pandoc invocation and are not
available to anything else.

A growing body of work documents the per-turn token cost of MCP servers
in agent workflows, with multi-server deployments routinely consuming
tens of thousands of tokens before useful work begins (Sadani and Kumar
2026; {``MCP {Servers} {Use} 35x {More} {Tokens} {Than} {CLI} {Tools}
--- {And} {Reliability} {Drops} to 72\% on {Hard} {Tasks}''} 2026).
Anthropic's analysis recommends presenting tools as code on a filesystem
and loading definitions on demand, reporting up to a 98.7\% reduction in
context overhead under that pattern ({``Code Execution with {MCP:}
Building More Efficient {AI} Agents \textbackslash{} {Anthropic},''}
n.d.). quartobot's design predates this literature but lands on the same
architecture: the CLI is the default surface, the MCP server is
optional, and the on-disk artifact is the structurally efficient locus
for downstream reuse.

\subsection{Reconciliation of resolved and hand-curated
entries}\label{reconciliation-of-resolved-and-hand-curated-entries}

Real manuscripts mix resolver-generated and hand-curated citations. A
\texttt{references.bib} may carry entries for textbooks, software
releases, gray literature, or any work without a persistent identifier
that the resolver covers. The \texttt{references.json} produced by
\texttt{quartobot\ resolve} is regenerated on each pre-render pass.
Without explicit handling, the two sources can collide on citation keys,
especially as resolver-generated keys are derived from the identifier
itself (e.g., \texttt{doi:10.1371/journal.pcbi.1007128}) and may overlap
conceptually with hand-curated keys for the same work.

The \texttt{reconcile} subcommand handles this case explicitly. Three
modes are supported: \texttt{prefer-bib} keeps the hand-curated entry
and discards the resolver's version; \texttt{prefer-resolved} keeps the
resolver's entry and removes the hand-curated duplicate; and
\texttt{fail-on-collision} halts with a diagnostic listing the
conflicting keys for the author to resolve. All three modes operate with
backup-then-mutate semantics: the original files are saved before any
change, so a \texttt{reconcile} step is reversible.

\section{Operation}\label{operation}

\subsection{Installation}\label{installation}

quartobot is distributed on PyPI and is installed as a standalone
command-line tool using \texttt{uv}:

\begin{Shaded}
\begin{Highlighting}[]
\ExtensionTok{uv}\NormalTok{ tool install quartobot}
\end{Highlighting}
\end{Shaded}

The \texttt{uv\ tool\ install} form is recommended over
\texttt{pip\ install} because it isolates quartobot in its own
environment and places the resulting executable on the user's
\texttt{PATH} without polluting any project's Python environment.
quartobot requires Python 3.10 or later; no other system dependencies
are required beyond an installed Quarto runtime.

\subsection{Canonical workflow}\label{canonical-workflow}

The shortest path from ``I want to write a manuscript'' to a working
manuscript-as-software project is three commands:

\begin{Shaded}
\begin{Highlighting}[]
\ExtensionTok{quarto}\NormalTok{ create project manuscript my{-}paper}
\BuiltInTok{cd}\NormalTok{ my{-}paper}
\ExtensionTok{quartobot}\NormalTok{ init}
\end{Highlighting}
\end{Shaded}

\texttt{quarto\ create} scaffolds a default Quarto manuscript with an
\texttt{index.qmd} source file, a placeholder \texttt{references.bib},
and a \texttt{\_quarto.yml} project configuration.
\texttt{quartobot\ init} then layers the citation pipeline on top: it
adds the \texttt{pre-render:} line to \texttt{\_quarto.yml}, ensures
\texttt{references.json} is listed in the \texttt{bibliography:} field,
and appends the relevant entries to \texttt{.gitignore} for any
transient files. The result is a Quarto project that resolves
\texttt{@doi:}, \texttt{@pmid:}, and other persistent-identifier
citation keys on every render.

Authors then write prose as they would in any Quarto project, with
citation keys pasted directly into the text:

\begin{Shaded}
\begin{Highlighting}[]
\NormalTok{We follow @doi:10.1371/journal.pcbi.1007128, with the dataset described}
\NormalTok{in @pmid:31479462 and methods inspired by @arxiv:2104.10729.}
\end{Highlighting}
\end{Shaded}

A \texttt{quarto\ render} runs the pre-render hook automatically,
produces an updated \texttt{references.json}, and proceeds with the
normal multi-format render.

To add the continuous-integration scaffolding described above, a single
additional command is run from inside the project directory:

\begin{Shaded}
\begin{Highlighting}[]
\ExtensionTok{quartobot}\NormalTok{ use github{-}ci}
\end{Highlighting}
\end{Shaded}

The author commits the generated workflow files and pushes; the first
push to GitHub initiates the render, publishes the result to GitHub
Pages, and registers the pull-request preview pattern for subsequent
contributions.

\subsection{Alternative on-ramps}\label{alternative-on-ramps}

The canonical workflow assumes a new project. Two alternative paths
cover other starting points.

For an \textbf{existing Quarto project} that already has prose and
possibly a \texttt{references.bib}, \texttt{quartobot\ init} is run from
inside the project root. It performs the same scaffolding steps and does
not modify existing prose or bibliography entries. Authors can then
begin adding \texttt{@doi:}-style keys incrementally without rewriting
any existing citations.

For the \textbf{resolver-only case} --- an author who wants citation
resolution but has their own CI configuration and does not want
quartobot to scaffold workflows --- \texttt{quartobot\ init} provides
everything needed without \texttt{quartobot\ use\ github-ci}. The author
wires the resolver into any existing pipeline by ensuring
\texttt{quartobot\ resolve} runs before pandoc.

\subsection{Command-line reference}\label{command-line-reference}

The full CLI is summarized in Table 2. Each subcommand has detailed help
available via
\texttt{quartobot\ \textless{}subcommand\textgreater{}\ -\/-help}.

{\def\LTcaptype{none} % do not increment counter
\begin{longtable}[]{@{}
  >{\raggedright\arraybackslash}p{(\linewidth - 2\tabcolsep) * \real{0.5000}}
  >{\raggedright\arraybackslash}p{(\linewidth - 2\tabcolsep) * \real{0.5000}}@{}}
\toprule\noalign{}
\begin{minipage}[b]{\linewidth}\raggedright
Subcommand
\end{minipage} & \begin{minipage}[b]{\linewidth}\raggedright
Purpose
\end{minipage} \\
\midrule\noalign{}
\endhead
\bottomrule\noalign{}
\endlastfoot
\texttt{resolve} & Scan project sources for citation keys, resolve each
via \texttt{citekey\_to\_csl\_item}, write CSL-JSON to
\texttt{references.json}. Invoked by Quarto's pre-render hook; also
directly callable. \\
\texttt{init} & Scaffold the citation pipeline into a Quarto project.
Idempotent. \\
\texttt{use\ github-ci} & Layer GitHub Actions workflows for render,
pull-request preview, and (optionally) per-commit permalinks. \\
\texttt{scan} & Inventory citation keys across a project without
resolving them. CI-lint surface. \\
\texttt{validate} & Static checks on \texttt{\_quarto.yml} and the
project layout. CI-lint surface. \\
\texttt{reconcile} & Resolve collisions between \texttt{references.bib}
and \texttt{references.json} with explicit modes (\texttt{prefer-bib},
\texttt{prefer-resolved}, \texttt{fail-on-collision}). \\
\texttt{versions} & Generate the \texttt{/versions/} index page on
\texttt{gh-pages} listing tagged releases and open pull-request
previews. \\
\texttt{mcp} & Run a stdio Model Context Protocol server exposing the
resolver for in-band use by AI authoring agents. \\
\end{longtable}
}

\section{Use cases}\label{use-cases}

\section{Discussion}\label{discussion}

\section{Software availability}\label{software-availability}

\begin{itemize}
\tightlist
\item
  \textbf{Source repository:}
  \href{https://github.com/seandavi/quartobot}{github.com/seandavi/quartobot}
\item
  \textbf{Package index:}
  \href{https://pypi.org/project/quartobot/}{\texttt{quartobot} on PyPI}
\item
  \textbf{Documentation:}
  \href{https://seandavi.github.io/quartobot/}{seandavi.github.io/quartobot}
\item
  \textbf{License:} MIT
\item
  \textbf{Archived release:} Zenodo DOI to be minted at submission.
\item
  \textbf{Latest version at submission:}
\end{itemize}

\section{Data availability}\label{data-availability}

This article describes a software tool; no primary research data are
generated. The two production case studies referenced in the Use Cases
section have their own data availability statements in their respective
publications. The reproducibility artifacts for those case studies ---
the source \texttt{.qmd} files, the committed \texttt{references.json}
bibliographies, the rendered outputs, and the GitHub Actions workflow
history --- are publicly available in the repositories cited under each
case study. The present manuscript is itself a quartobot project; its
source, including the committed \texttt{references.json} produced by
\texttt{quartobot\ resolve}, is available at .

\section{Author contributions}\label{author-contributions}

\begin{itemize}
\tightlist
\item
  \textbf{Sean Davis} --- Conceptualization, Methodology, Software,
  Investigation, Writing --- Original Draft, Writing --- Review \&
  Editing, Project Administration.
\end{itemize}

\section{AI usage statement}\label{ai-usage-statement}

Portions of this manuscript were drafted with the assistance of large
language model coding agents (Anthropic Claude, operating through Claude
Code) used as a literate-programming and drafting tool. AI assistance
was used during outlining, prior-work surveying, structural drafting of
the Implementation and Operation sections, and the development of the
Discussion. All scientific claims, design decisions, citation choices,
and final wording were reviewed and edited by the human authors.
quartobot itself was developed by the human authors with AI assistance
in implementation and documentation.

\section{Funding}\label{funding}

This work was supported by the National Cancer Institute under the NCI
Information Technology for Cancer Research award U24CA289073 (NCI
Bioconductor) and the University of Colorado Cancer Center Support Grant
P30CA046934.

\section{Acknowledgments}\label{acknowledgments}

The authors thank the manubot project and its maintainers for eight
years of resolver expertise embodied in \texttt{citekey\_to\_csl\_item},
on which quartobot depends; the Quarto team at Posit for the Quarto
Manuscripts project type that made this integration tractable; and early
users of quartobot in the Bioconductor and CFDE communities whose
feedback shaped the tool.

\section{References}\label{references}

\protect\phantomsection\label{refs}
\begin{CSLReferences}{1}{1}
\bibitem[\citeproctext]{ref-url:https:ux2fux2fwww.anthropic.comux2fengineeringux2fcode-execution-with-mcp}
{``Code Execution with {MCP:} Building More Efficient {AI} Agents
\textbackslash{} {Anthropic}.''} n.d. Accessed May 28, 2026.
\url{https://www.anthropic.com/engineering/code-execution-with-mcp}.

\bibitem[\citeproctext]{ref-url:https:ux2fux2fwww.mindstudio.aiux2fblogux2fmcp-servers-35x-more-tokens-cli-tools-reliability-benchmark}
{``MCP {Servers} {Use} 35x {More} {Tokens} {Than} {CLI} {Tools} ---
{And} {Reliability} {Drops} to 72\% on {Hard} {Tasks}.''} 2026. May 9.
\url{https://www.mindstudio.ai/blog/mcp-servers-35x-more-tokens-cli-tools-reliability-benchmark/}.

\bibitem[\citeproctext]{ref-arxiv:2604.21816}
Sadani, Anuj, and Deepak Kumar. 2026. \emph{Tool {Attention} {Is} {All}
{You} {Need:} {Dynamic} {Tool} {Gating} and {Lazy} {Schema} {Loading}
for {Eliminating} the {MCP/Tools} {Tax} in {Scalable} {Agentic}
{Workflows}}. 2604.21816. arXiv. \url{https://arxiv.org/abs/2604.21816}.

\end{CSLReferences}




\end{document}
